\section{客观要件 - 结果}
\begin{enumerate}
  \item 分量不是那么重。
  \item 法益侵害事实分类。
        \begin{enumerate}
          \item 实害犯。成立这个罪要具备这个实害结果。\al{生产销售劣药罪，吃死人}。
          \item 危险犯。
                \begin{enumerate}
                  \item 具体危险犯。实施 + 产生紧迫的危险。
                  \item 抽象危险犯/行为犯。只要实施了，就构成。
                        \begin{enumerate}
                          \item 危险程度很遥远，很缓和。
                          \item \al{生产、销售、提供假药罪，不要求吃死人 \hz 危害身体健康}。
                        \end{enumerate}
                \end{enumerate}
          \item \xjt 含义很混乱，考的少，了解就行。
          \item 行为犯。有行为即可，不要求实害结果。\al{伪证罪，不需要让人真的坐牢才成立}。
          \item 结果犯。（严重）实害结果作为犯罪成立的条件。\al{丢失枪支 \yjt 捡到交给警察、滥用职权 \yjt 重大损失}。
        \end{enumerate}
  \item \zd{结果加重犯}。其实就是想象竞合犯（\zd{一个行为}），但法条做了特殊规定（\zd{法定性}）。
        \begin{enumerate}
          \item 最典型。故意伤害罪（\zd{基本犯}）致人死亡（\zd{加重结果}）。
          \item 基本犯都是故意，加重结果基本是过失。例外。加重结果过失 \hz 故意均可。\zd{强奸、拐卖、抢劫、放火致人重伤死亡}。
          \item 因果关系。\zd{加重结果必须是基本犯的实行行为导致的}。
                \begin{enumerate}
                  \item 判断方法。\zd{行为与目的同时存在原则}。一个行为可以成为 N 多罪的实行行为，看带着啥目的实施的实行行为。\zd{不看既遂未遂}。
                  \item \al{甲抢劫暴力致乙死亡} \yjt 取财目的。
                  \item \al{甲抢劫发现对方是仇人杀人} \yjt 杀人目的。
                  \item \al{甲强奸乙女被乙女咬断舌头后泄愤杀人} \yjt 杀人目的。
                  \item \al{甲抢劫杀人毁尸灭迹}。
                  \item \al{甲抢劫乙，枪头砸乙背部，不慎走火杀死了乙}。
                \end{enumerate}
        \end{enumerate}
\end{enumerate}
